\documentclass{beamer}
\usetheme[style=beta]{epyt}
\usepackage{macros_beamer}
\usepackage[portuguese]{babel}
\title{Introdução à Teoria Espectral de Grafos}
\subtitle{Passeios aleatórios em grafos}
\date{03 de julho de 2026}
\author{Leonardo Suarez}
\institute{Instituto de Matemática e Estatística - UFRGS}

\usepackage[backend=biber, style=ieee]{biblatex}\addbibresource{referencias_grafos.bib}
\setbeamertemplate{bibliography item}{\insertbiblabel}

\DeclareMathOperator{\tr}{tr}
\DeclareMathOperator{\spect}{spect}

\begin{document}

\nocite{spielman}

\frame{\titlepage}

\begin{frame}
    \tableofcontents
\end{frame}

\section{Passeios aleatórios}

\begin{frame}{definições}
    Seja \(G = (V, E, w)\) um grafo valorado não direcionado.
    
    \begin{itemize}
        \item um \alert{passeio aleatório} em \(G\) é um processo que começa num vértice e a cada instante se move para um vértice vizinho;
        \item quando \(G\) é valorado, o passeio se move com probabilidade proporcional ao peso das arestas.
    \end{itemize}

    Denotamos por \(\mathbf{p}_t \in \R^n\) a distribuição de probabilidade no instante \(t\) e escrevemos \(\mathbf{p}_t(a)\) para indicar a probabilidade de estar no vértice \(a \in V\) no instante \(t\). 

    Um vetor probabilidade \(\mathbf{p}\) é um vetor tal que \(\mathbf{p}(a) \geq 0\) para todo \(a \in V\) e
    \begin{equation*}
        \sum_{a} \mathbf{p}(a) = 1.
    \end{equation*}
\end{frame}

\begin{frame}{matriz de transição}
    Observamos que a distribuição de probabilidade no instante \(t+1\) é dada por
    \begin{equation}
        \mathbf{p}_{t+1} = \sum_{b \mid b \sim a} \frac{w(a, b)}{\mathbf{d}(b)} \mathbf{p}_t(b),
    \end{equation} onde \(\mathbf{d}(b) = \sum_{a} w(a, b)\) é o grau do vértice \(b\). Em notação matricial, definimos a \alert{matriz de transição} \(\mathbf{W}\)

    \begin{equation}
        \mathbf{W} := \mathbf{M} \mathbf{D}^{-1},
    \end{equation} de modo que \(\mathbf{p}_{t+1} = \mathbf{W} \mathbf{p}_t = \mathbf{W}^{t+1}\mathbf{p}_0\).
\end{frame}

\begin{frame}{matriz de transição preguiçosa}
    Frequentemente consideramos \alert{passeios aleatórios preguiçosos}, que evoluem conforme

    \begin{equation}
        \mathbf{p}_{t+1} = \frac{1}{2}\mathbf{p}_t(a) + \frac{1}{2}\sum_{b \mid b \sim a} \frac{w(a, b)}{\mathbf{d}(b)} \mathbf{p}_t(b),
    \end{equation} e satisfazem
    \begin{equation}
        \mathbf{p}_{t+1} = \widetilde{\mathbf{W}} \mathbf{p}_t,
    \end{equation} onde \(\widetilde{\mathbf{W}} := \frac{1}{2}\mathbf{I} + \frac{1}{2}\mathbf{W}\) é a \alert{matriz de transição preguiçosa}.
\end{frame}

\section{Espectro de matrizes de transição}

\begin{frame}{o Teorema de Perron-Frobenius}
    \begin{thm}[Perron-Frobenius, caso simétrico]
        Sejam \(G\) um grafo conexo valorado, \(\mathbf{M}\) a sua matriz de adjacência, \(\mathbf{H}\) uma matriz diagonal não negativa, \(\mathbf{M}^{\prime} = \mathbf{M} + \mathbf{H}\) e \(\mu_1 \geq \mu_2 \geq \ldots \geq \mu_n\) os autovalores de \(\mathbf{M}^{\prime}\). Então
        \begin{enumerate}
            \item o autovalor \(\mu_1\) tem um autovetor com entradas estritamente positivas;
            \item \(\mu_1 \geq -\mu_n\);
            \item \(\mu_1 > \mu_2\), i.e., o autovalor \(\mu_1\) é simples;
        \end{enumerate}
    \end{thm}
\end{frame}

\begin{frame}{matriz de adjacência normalizada}
    Definimos a \alert{matriz de adjacência normalizada} de \(G\) como
    \begin{equation}
        \mathbf{A} := \mathbf{D}^{-1/2} \mathbf{W} \mathbf{D}^{1/2} = \mathbf{D}^{-1/2} \mathbf{M} \mathbf{D}^{-1/2},
    \end{equation} de modo que \(\mathbf{A}\) é simétrica.

    \begin{prop}
        O vetor \(\mathbf{\psi}\) é autovetor de \(\mathbf{A}\) com autovalor \(\omega\) se, e somente se, \(\mathbf{D}^{-1/2}\mathbf{\psi}\) é autovetor de \(\mathbf{W}\) com autovalor \(\omega\).
    \end{prop}
\end{frame}

\begin{frame}{consequências de Perron-Frobenius}
    \begin{itemize}
        \item \(\widetilde{\mathbf{W}}\) tem os mesmos autovetores de \(\mathbf{W}\);
        \item o vetor dos graus \(\mathbf{d}\) é \alert{vetor de Perron} de \(\mathbf{W}\) com autovalor \(1\) (por quê?);
        \item \(\mathbf{\psi}_1 := \frac{\mathbf{d}^{1/2}}{\Vert \mathbf{d}^{1/2} \Vert}\) é o vetor de Perron unitário de \(\mathbf{A}\) com autovalor \(1\);
        \item o Teorema de Perron-Frobenius garante que todos os autovalores de \(\mathbf{A}\) estão entre \(-1\) e \(1\);
        \item como \(\widetilde{\mathbf{W}} = \frac{1}{2}\mathbf{I} + \frac{1}{2}\mathbf{W}\), todos os autovalores de \(\widetilde{\mathbf{W}}\) estão entre \(0\) e \(1\).
    \end{itemize}
\end{frame}

\section{A distribuição estacionária}

\begin{frame}
    
\end{frame}

\begin{frame}{Taxa de convergência}
    
\end{frame}

\section{Relação com a matriz laplaciana normalizada}

\begin{frame}{Matriz laplaciana normalizada}
    
\end{frame}

\begin{frame}
    
\end{frame}

\begin{frame}
    
\end{frame}

\section{Exemplos}


\section*{Referências}

\begin{frame}
    \printbibliography
\end{frame}

\end{document}
